% Options for packages loaded elsewhere
\PassOptionsToPackage{unicode}{hyperref}
\PassOptionsToPackage{hyphens}{url}
\documentclass[
  ignorenonframetext,
]{beamer}
\newif\ifbibliography
\usepackage{pgfpages}
\setbeamertemplate{caption}[numbered]
\setbeamertemplate{caption label separator}{: }
\setbeamercolor{caption name}{fg=normal text.fg}
\beamertemplatenavigationsymbolsempty
% remove section numbering
\setbeamertemplate{part page}{
  \centering
  \begin{beamercolorbox}[sep=16pt,center]{part title}
    \usebeamerfont{part title}\insertpart\par
  \end{beamercolorbox}
}
\setbeamertemplate{section page}{
  \centering
  \begin{beamercolorbox}[sep=12pt,center]{section title}
    \usebeamerfont{section title}\insertsection\par
  \end{beamercolorbox}
}
\setbeamertemplate{subsection page}{
  \centering
  \begin{beamercolorbox}[sep=8pt,center]{subsection title}
    \usebeamerfont{subsection title}\insertsubsection\par
  \end{beamercolorbox}
}
% Prevent slide breaks in the middle of a paragraph
\widowpenalties 1 10000
\raggedbottom
\AtBeginPart{
  \frame{\partpage}
}
\AtBeginSection{
  \ifbibliography
  \else
    \frame{\sectionpage}
  \fi
}
\AtBeginSubsection{
  \frame{\subsectionpage}
}
\usepackage{iftex}
\ifPDFTeX
  \usepackage[T1]{fontenc}
  \usepackage[utf8]{inputenc}
  \usepackage{textcomp} % provide euro and other symbols
\else % if luatex or xetex
  \usepackage{unicode-math} % this also loads fontspec
  \defaultfontfeatures{Scale=MatchLowercase}
  \defaultfontfeatures[\rmfamily]{Ligatures=TeX,Scale=1}
\fi
\usepackage{lmodern}
\ifPDFTeX\else
  % xetex/luatex font selection
\fi
% Use upquote if available, for straight quotes in verbatim environments
\IfFileExists{upquote.sty}{\usepackage{upquote}}{}
\IfFileExists{microtype.sty}{% use microtype if available
  \usepackage[]{microtype}
  \UseMicrotypeSet[protrusion]{basicmath} % disable protrusion for tt fonts
}{}
\makeatletter
\@ifundefined{KOMAClassName}{% if non-KOMA class
  \IfFileExists{parskip.sty}{%
    \usepackage{parskip}
  }{% else
    \setlength{\parindent}{0pt}
    \setlength{\parskip}{6pt plus 2pt minus 1pt}}
}{% if KOMA class
  \KOMAoptions{parskip=half}}
\makeatother
\setlength{\emergencystretch}{3em} % prevent overfull lines
\providecommand{\tightlist}{%
  \setlength{\itemsep}{0pt}\setlength{\parskip}{0pt}}
% Auto-generated by infrastructure/rendering/_pdf_latex_helpers.py::extract_math_font_preamble.
% Minimal header passed to Pandoc via -H so Beamer slides pick up the same math font
% as the combined-PDF path. Adjust the manuscript's preamble.md to change the font globally.
\usepackage{unicode-math}
\setmathfont{latinmodern-math.otf}

% Slide layout defaults for warning-clean scientific decks.
\usepackage{etoolbox}
\IfFileExists{xurl.sty}{\usepackage{xurl}}{}
\IfFileExists{seqsplit.sty}{\usepackage{seqsplit}}{\newcommand{\seqsplit}[1]{#1}}
\protected\def\breakseq#1{\seqsplit{#1}}
\protected\def\breaktt#1{\begingroup\ttfamily\seqsplit{#1}\endgroup}
\setlength{\emergencystretch}{6em}
\tolerance=5000
\hbadness=10000
\hfuzz=1pt
\setlength{\tabcolsep}{2pt}
\AtBeginEnvironment{longtable}{\tiny\renewcommand{\arraystretch}{0.86}\setlength{\tabcolsep}{1pt}}
\AtBeginEnvironment{tabular}{\tiny\renewcommand{\arraystretch}{0.86}\setlength{\tabcolsep}{1pt}}

% Natbib and cross-reference fallbacks — slides don't load natbib
% or cleveref, but combined-PDF manuscript prose may emit these
% commands. The fallback renders citations as a bracketed key list
% and cross-references as detokenized labels so slides don't fail on
% undefined control sequences or raw underscores. \providecommand is
% a no-op if packages load later.
\providecommand{\citep}[1]{[#1]}
\providecommand{\citet}[1]{#1}
\providecommand{\citealp}[1]{#1}
\providecommand{\citeauthor}[1]{#1}
\providecommand{\citeyear}[1]{#1}
\providecommand{\cref}[1]{\texttt{\detokenize{#1}}}
\providecommand{\Cref}[1]{\texttt{\detokenize{#1}}}
\usepackage{bookmark}
\IfFileExists{xurl.sty}{\usepackage{xurl}}{} % add URL line breaks if available
\urlstyle{same}
% fallback for those not using the hyperref driver hyperxmp:
\makeatletter
\@ifundefined{xmpquote}{\newcommand{\xmpquote}[1]{#1}}{}
\makeatother
\hypersetup{
  hidelinks,
  pdfcreator={LaTeX via pandoc}}

\author{\texorpdfstring{}{}}
\date{}

\begin{document}

\section{Scope, Related Work, and Positioning}\label{sec:scope}

\begin{frame}[allowframebreaks]{Scope, Related Work, and Positioning}
This section situates the exemplar and states explicit boundaries. The
goal is not to compete with comprehensive treatments of exploratory data
analysis \breakseq{[@tukey1977eda]} or statistical graphics
\breakseq{[@wilkinson2005grammar]}, but to show how a minimal, test-backed EDA
story fits the template's reproducibility and rendering stack
\breakseq{[@peng2011reproducible]}.
\end{frame}

\begin{frame}[allowframebreaks]{Exploratory data analysis}
\protect\phantomsection\label{exploratory-data-analysis}
The practice of examining a dataset before formal modelling --- looking
at distributions, missingness, and relationships --- traces to Tukey's
foundational work \breakseq{[@tukey1977eda]} and is supported in modern
practice by the pandas data analysis toolkit \breakseq{[@mckinney2010pandas]}
and the broader scientific-Python stack \breakseq{[@harris2020numpy]}. The
present manuscript restricts attention to a \textbf{first pass} on a
\textbf{small tabular dataset}: load, surface missingness, compute
descriptive statistics and per-group means, and rank features by Pearson
correlation.
\end{frame}

\begin{frame}[allowframebreaks]{Modelling and inference (out of scope)}
\protect\phantomsection\label{modelling-and-inference-out-of-scope}
What comes \emph{after} a first EDA pass --- hypothesis testing,
regression, dimensionality reduction, or predictive modelling --- is
deliberately \textbf{out of scope}. The exemplar keeps the analysis
minimal so the architectural lesson (notebook -\textgreater{} tested src
extraction) stays visible rather than buried under statistical
machinery.
\end{frame}

\begin{frame}[fragile,allowframebreaks]{What this project proves about
the template}
\protect\phantomsection\label{what-this-project-proves-about-the-template}
The analytical steps here are standard. The \textbf{non-standard}
contribution is procedural: the same tested functions in
\texttt{src/eda/} back the interactive notebook, the headless analysis
script, and this manuscript, so the figures and the summary table always
refer to the same code. That pattern is what downstream projects should
copy --- whether the domain is survey data, sensor logs, or experimental
measurements.
\end{frame}

\begin{frame}[allowframebreaks]{Explicit limitations}
\protect\phantomsection\label{explicit-limitations}
\begin{enumerate}
\tightlist
\item
  \textbf{Dataset size}: a single small synthetic cohort (120 rows) is
  used for transparent, fast reproduction; no large-scale or streaming
  data is handled.
\item
  \textbf{Cleaning policy}: only listwise deletion is implemented;
  imputation is left as a documented extension.
\item
  \textbf{Correlation method}: only Pearson correlation is computed;
  rank-based (Spearman) or non-linear association measures are out of
  scope.
\item
  \textbf{Statics only}: the library returns plot-ready data;
  interactive widgets and dashboards are not part of the exemplar.
\end{enumerate}

These limitations are intentional: they narrow the surface so that the
reproducibility concerns --- tested functions, a thin script, and a
structurally verified notebook --- remain visible rather than buried
under analytical complexity.
\end{frame}

\end{document}
